\section{Multi-Modal Assessment Framework}

\label{sec:assessment}

\subsection{Assessment Modalities}

ZooEd employs four assessment modalities, each calibrated for reliability across languages and cultural contexts:

\begin{table}[H]
\centering
\caption{Assessment modalities and their characteristics.}
\label{tab:assessment_modalities}
\begin{tabularx}{\textwidth}{@{}lccX@{}}
\toprule
\textbf{Modality} & \textbf{Objectives} & \textbf{Reliability} & \textbf{Description} \\
\midrule
Image Identification & 834 & $\kappa = 0.91$ & Species identification from photographs with graduated difficulty \\
Scenario Reasoning & 1,176 & $\kappa = 0.84$ & Multi-step decision scenarios with branching consequences \\
Free Response & 612 & $\kappa = 0.79$ & Open-ended explanations evaluated by calibrated LLM rubrics \\
Practical Demonstration & 225 & $\kappa = 0.88$ & Video or photo evidence of field skill execution \\
\bottomrule
\end{tabularx}
\end{table}

\subsection{Image-Based Species Identification}

Species identification assessments use a graduated difficulty framework.
For each target species $s$, we construct distractor sets at five difficulty levels:

\begin{equation}
    D_k(s) = \{s' \in \mathcal{S} : d_{\text{taxon}}(s, s') \leq k \wedge d_{\text{visual}}(s, s') \leq \epsilon_k\}
\end{equation}

\noindent where $d_{\text{taxon}}$ is taxonomic distance (same genus = 1, same family = 2, etc.), $d_{\text{visual}}$ is visual similarity computed from a pretrained feature extractor, and $\epsilon_k$ controls the visual confusability at level $k$.

Assessment images are drawn from the Zoo Data Commons (\Cref{sec:data_commons_ref}), filtered for quality and annotated with difficulty metadata.
Each assessment presents 4--6 candidate species with the learner selecting the correct identification and optionally providing a justification that references diagnostic features.

\subsection{Scenario-Based Reasoning}

Scenario assessments present multi-step conservation decision problems.
Each scenario is modeled as a decision tree:

\begin{equation}
    \mathcal{T} = (V, E, \rho, \omega)
\end{equation}

\noindent where $V$ is the set of decision nodes, $E$ is the set of directed edges (choices), $\rho: V \rightarrow \mathcal{O}^*$ maps each node to the set of objectives it assesses, and $\omega: E \rightarrow \mathbb{R}$ assigns ecological outcome scores to each decision path.

The system evaluates not only the terminal outcome but the reasoning trajectory, awarding partial credit for locally optimal decisions even when the global path is suboptimal.
This rubric is computed as:

\begin{equation}
    \text{Score}(\pi) = \sum_{v \in \pi} w_v \cdot \mathbb{1}[\text{choice}(v) \in \text{acceptable}(v)]
\end{equation}

\noindent where $\pi$ is the learner's path through the scenario, $w_v$ is the weight of node $v$ (proportional to the difficulty of the assessed objectives), and $\text{acceptable}(v)$ is the set of defensible choices at node $v$ as determined by expert annotation.

\subsection{LLM-Based Free Response Evaluation}

Free-response answers are evaluated using a calibrated language model pipeline:

\begin{enumerate}[leftmargin=*]
    \item \textbf{Rubric Application:} The LLM evaluates the response against a structured rubric specific to the learning objective, producing scores on 3--5 criteria.
    \item \textbf{Multi-Rater Simulation:} Three independent LLM evaluations with different random seeds produce a score distribution, mitigating single-evaluation bias.
    \item \textbf{Calibration:} Scores are calibrated against a held-out set of 500 expert-graded responses per language, using isotonic regression to map raw LLM scores to calibrated mastery probabilities.
    \item \textbf{Feedback Generation:} The LLM generates formative feedback in the learner's language, identifying strengths and specific areas for improvement.
\end{enumerate}

The calibrated mastery estimate from free response assessment is:

\begin{equation}
    \hat{p}_{\text{mastery}} = f_{\text{cal}}\left(\frac{1}{K}\sum_{k=1}^{K} r_k\right)
\end{equation}

\noindent where $r_k$ is the raw score from the $k$-th LLM evaluation and $f_{\text{cal}}$ is the calibration function.

% ============================================================
% 6. System Architecture
% ============================================================
